\begin{python0}
from solutions import *; clear()
\end{python0}
\sectionthree{de Morgan's laws}

Look at this statement. John said to Mary:

\[\text{“I'm not stupid and I'm not ugly!!!”}\]

What about this. John said

\[\text{“It's not true that I'm either stupid or ugly!!!”}\]

Note that they both mean the same thing. In other words if b and c are
boolean values, then

\[\text{(!b) \&\& (!c)}\]

is the same as

\[\text{!(b $||$ c)}\]

The two facts:

\begin{tabular}{p{3cm} p{3cm} p{3cm}}
(!b) \&\& (!c) & is the same as & !(b \&\& c)\\
(!b) $||$ (!c) & is the same as & !(b $||$ c)\\
\end{tabular}

are called \EMPHASIZE{de Morgan's laws.}

Why are these formula helpful? Here's one use of de Morgan's Laws to
simplify a boolean expression:

\[\text{!(x $<$ 0 $||$ x $>$ 5) = !(x $<$ 0) \&\& !(x $>$ 5)}\]

Now note that !(x $<$ 0) is the same as x $>=$ 0. Right? Also, !(x $>$ 5) is the same as x $<=$ 5. So

\[\text{!(x $<$ 0 $||$ x $>$ 5) =! (x $<$ 0) \&\& !(x $>$ 5)}\]
\[ \tab[3em]{\text{= (x $>=$ 0) \&\& (x $<=$ 5)}}\]

Note that the final boolean expression does not have the ! operator.
Altogether

\begin{tabular}{p{5cm} p{5cm}}
!(x $<$ 0 $||$ x $>$ 5) & uses 4 operators\\
(x $>=$ 0) \&\& (x $<=$ 5) & uses 3 operators\\
\end{tabular}

Get it?

\textbf{Exercise.} Using de Morgan's law, simplify this expression:

\[\text{!(x $<$ y + z \&\& y – z == x)}\]

\textbf{Exercise.} Using de Morgan's law, simplify this expression:

\[\text{!(x $<$ y + z $||$ y != x + y \&\& y – z == x)}\]

\textbf{Exercise.} Using de Morgan's law, simplify this expression:

\[\text{!((x < y + z $||$ y != x + y) \&\& y – z == x)}\]

Write a simple C++ program to verify your simplification.

\textbf{Exercise.} Simplify this expression:

\[\text{!((x $<$ y * z $||$ y \% 3 $<$ 1) \&\& y $>=$ x + 1)}\]

(All variables are integer variables with positive values.) Write a simple
C++ program to verify your simplification.

